\steppingstone{Teaching a Machine}
\label{stone:computer}

By the time we reach this point, the arithmetic is finished. We know how to measure, how to grow a construction, how to multiply and divide exactly, how to move through space, and how to reflect a structure without losing its exactness. The remaining question is no longer mathematical. It is practical.

Can a computer use this arithmetic just as faithfully as we do?

A computer does not understand mathematics. It does not recognize elegance, nor does it know whether one method is better than another. It simply performs the arithmetic it has been given, over and over again, without becoming tired or distracted. If the arithmetic is good, the computation is good. If the arithmetic is poor, the computer faithfully repeats the same poor method millions of times.

That observation shifts our attention away from the machine and back toward the arithmetic itself. Every computer program eventually becomes an enormous collection of very small operations. Numbers are added, subtracted, multiplied, divided, stored, and retrieved. The larger computation is nothing more than those small steps repeated many times.

Builders learned long ago that repeated work should be avoided whenever possible. A cabinet maker keeps templates. A machinist keeps jigs. An architect keeps drawings. The successful construction is preserved so that it never has to be rediscovered. Throughout this book the Fibonacci Spine, the Dense Grid, and the PhiBase Slide Rule have all followed that same idea. Do the difficult work once. Remember it. Reuse it whenever it is needed.

A computer is particularly good at remembering. Once an exact PhiBase relationship has been constructed, the machine can retrieve it whenever the same situation appears again. Nothing about the mathematics changes. Only the amount of repeated work changes.

This is not a new engineering idea. One of the most successful early microprocessors, the MOS 6502, contained no hardware multiplication instruction. Programmers frequently replaced repeated multiplication with carefully prepared lookup tables because memory was often less expensive than arithmetic. PhiBase follows exactly the same philosophy. The mathematics is performed once. The result is remembered. The computer simply repeats the construction.

Eventually every mathematical idea reaches the same moment. Someone builds it.

To answer the obvious question, an existing transformer inference engine was rebuilt so that its arithmetic used PhiBase instead of conventional floating-point multiplication. The model itself did not change. The training did not change. Only the arithmetic changed. The same inputs were presented, the same outputs were compared, and the resulting embeddings closely matched those of the original implementation while requiring substantially fewer arithmetic cycles.

That experiment does not prove that every future computer should use PhiBase. It answers a simpler question. Can a modern machine perform exact PhiBase arithmetic throughout a practical computation?

Yes.

Artificial intelligence is only one example. Robotic construction, structural analysis, crystal growth, and other problems built upon golden-ratio geometry are equally natural applications. Wherever the geometry belongs, the arithmetic can follow.

\section*{Builder's Notebook}

Imagine two workshops. In the first, every successful measurement is repeated from the beginning each morning. In the second, every successful construction is carefully recorded and reused. Both workshops produce the same object. Which one would you rather work in? Now replace the workshop with a computer. Nothing else has changed.

\section*{Looking Ahead}

We began with two simple triangles. We now have an arithmetic that both people and machines can follow exactly. One question remains.

What have we really been building?
